Per a identificar els instruments musicals es pot utilitzar un espectre de freqüències. A la figura següent es representa l'espectre d'un instrument que es vol identificar. L'oboè és un instrument que té un extrem obert i un extrem tancat que és l'embocadura amb una canya. A l'extrem obert, l'amplitud de la vibració de les molècules és màxima: tenim un ventre. En canvi, a l'altre extrem, el tudell per on es bufa és una canya que comunica pressió a l'aire i impedeix que aquest es pugui moure amb llibertat: és un node. D'altra banda, en un piano les cordes estan pinçades pels dos extrems, és a dir, els dos extrems de la corda són nodes.

\figura[0.6]{fig1.png}

\begin{apartat}{1,25}
Quina és la freqüència fonamental d'aquest espectre? Determineu si es tracta d'un oboè o d'un piano, i justifiqueu la resposta. Determineu també la llargària del tub o de la corda.
\end{apartat}

\begin{apartat}{1,25}
Un quintet de vent acostuma a estar format per cinc instrumentistes, normalment flauta travessera, oboè, clarinet, trompa i fagot. A l'inici de la interpretació d'una composició tots 5 toquen de forma suau i amb la mateixa intensitat, generant un nivell d'intensitat sonora de 80 dB a un espectador que es troba a una distància d'1,5 m. Quina és la potència sonora de cada instrument? Suposeu que aquesta potència es distribueix uniformement per tota l'àrea d'una semiesfera.
\end{apartat}

\dades{%
  $I_0 = 10^{-12}\ \mathrm{W\,m^{-2}}$.\\
  La velocitat del so en l'aire és de $340\ \mathrm{m\,s^{-1}}$.}
